\documentclass[11pt]{article}
\usepackage[a4paper,margin=1in]{geometry}
\usepackage{amsmath,amssymb}
\usepackage[T1]{fontenc}
\usepackage{lmodern}
\usepackage{microtype}
\usepackage{xcolor}

\setlength{\parindent}{0pt}
\setlength{\parskip}{0.5em}

\newcommand{\promptbox}[1]{%
  \begingroup\setlength{\fboxsep}{10pt}\noindent
  \colorbox{blue!8}{\parbox{\dimexpr\linewidth-2\fboxsep\relax}{\textbf{Prompt. }#1}}
  \endgroup
}
\newcommand{\resultbox}[1]{%
  \begingroup\setlength{\fboxsep}{10pt}\noindent
  \colorbox{purple!8}{\parbox{\dimexpr\linewidth-2\fboxsep\relax}{\textbf{Result. }#1}}
  \endgroup
}
\newcommand{\examplebox}[1]{%
  \begingroup\setlength{\fboxsep}{10pt}\noindent
  \colorbox{green!8}{\parbox{\dimexpr\linewidth-2\fboxsep\relax}{\textbf{Example. }#1}}
  \endgroup
}

\title{\textbf{Book 1 --- Lecture 1 --- Exercise 3}\\[2mm]
\large Allowable Laws on the Integer Line}
\author{}
\date{}

\begin{document}
\maketitle

\promptbox{Determine which of the dynamical laws shown in Eqs.\ (2) through (5)
are \emph{allowable}. Equation (1), \(N(n+1)=N(n)-1\), is given as an allowable
example.}

\section*{What ``allowable'' means}
A law is allowable when it is:
\begin{enumerate}
  \item \textbf{Deterministic}: each current state has one successor.\footnote{A
  \emph{successor} is the state reached in one forward update.}
  \item \textbf{Reversible}: each state has one predecessor\footnote{A
  \emph{predecessor} is a state that maps to the current one under the update
  rule.} (bijection\footnote{A \emph{bijection} is both one-to-one and onto:
  every target has exactly one source.} on \(\mathbb{Z}\)\footnote{\(\mathbb{Z}\)
  denotes all integers: \(\dots,-2,-1,0,1,2,\dots\).}).
\end{enumerate}

\section*{Equations}
\[
\begin{aligned}
(1)\quad &N(n+1)=N(n)-1,\\
(2)\quad &N(n+1)=N(n)+1,\\
(3)\quad &N(n+1)=N(n)+2,\\
(4)\quad &N(n+1)=N(n)^2,\\
(5)\quad &N(n+1)=(-1)^{N(n)}N(n).
\end{aligned}
\]

\section*{Verdicts}
\textbf{(2) \(N(n+1)=N(n)+1\): allowable.}\\
It is the inverse-time version of Eq.\ (1), with inverse \(N\mapsto N-1\).

\textbf{(3) \(N(n+1)=N(n)+2\): allowable.}\\
It is a bijection with inverse \(N\mapsto N-2\).

\textbf{(4) \(N(n+1)=N(n)^2\): not allowable.}\\
Not injective\footnote{\emph{Injective} means different inputs give different
outputs.} (\(2^2=(-2)^2\)) and not surjective\footnote{\emph{Surjective} means
every target value is hit by at least one input.} onto \(\mathbb{Z}\).

\textbf{(5) \(N(n+1)=(-1)^{N(n)}N(n)\): allowable.}\\
Even integers are fixed; odd integers map to their negative, giving 2-cycles.
Each state has a unique predecessor.

\section*{Summary}
\[
\begin{array}{c|c}
\text{Equation} & \text{Allowable?} \\
\hline
(2)\ N(n+1)=N(n)+1 & \text{Yes} \\
(3)\ N(n+1)=N(n)+2 & \text{Yes} \\
(4)\ N(n+1)=N(n)^2 & \text{No} \\
(5)\ N(n+1)=(-1)^{N(n)}N(n) & \text{Yes}
\end{array}
\]

\resultbox{Allowable: (2), (3), and (5). Not allowable: (4).}

\examplebox{Eq.\ (4) fails reversibility immediately because
\(2^2=(-2)^2=4\), so the predecessor of 4 is not unique.}

\end{document}
